\documentclass[11pt,a4paper]{article}
\usepackage{curso}
\begin{document}
\portada{Ejercicios --- Lección}{5.5}{Otros sustratos físicos}
\section{Ejercicios}
\begin{enumerate}[label=\textbf{E\arabic*.}]
\item (B) Tiempo medio de primer paso de una caminata libre sobre $n$ sitios con inicio reflectante ($k=1$).
\item (B) Entropía disipada por paso con tasa hacia adelante doble de la de retroceso.
\item (B) ¿Por qué un circuito browniano necesita trinquetes?
\item (B) Escriba la dinámica de una red RC y su estado estacionario.
\item (I) Con $k_+=2$, $k_-=1$ y $n=20$, ¿cuánta entropía total se disipa en promedio? Compare con la cota de la TUR para $\sigma/\langle T\rangle=0{,}38$.
\item (I) Para un ciclo de $r$ nodos, escriba los modos de $L$ y el tiempo de mezcla de la difusión RC.
\item (A) Simule un CJoin con dos tokens independientes y mida el tiempo de espera hasta que ambos llegan.
\end{enumerate}
\section{Soluciones}
\begin{solucion}
\textbf{E1.} $\langle T\rangle=n(n+1)/2$ (para $n=20$: 210).

\textbf{E2.} $\ln2\approx0{,}69\,\kB$.

\textbf{E3.} Sin ellos la difusión no tiene dirección: el tiempo crece como $n^2$ y el cómputo puede deshacerse.

\textbf{E4.} $C\dot V=-LV$; estacionario: voltaje uniforme igual al promedio inicial (si la red es conexa).

\textbf{E5.} $\Sigma\approx n\ln2=13{,}9$; TUR: $\epsilon^2\ge2/\Sigma$ da $0{,}144$, y $0{,}38^2=0{,}144$: prácticamente saturada.

\textbf{E6.} Autovalores $2-2\cos(2\pi k/r)$; tiempo de mezcla $\sim C/(2-2\cos(2\pi/r))\approx Cr^2/(4\pi^2)$.

\textbf{E7.} El tiempo es el máximo de dos primeros pasos: mayor media que uno solo.
\end{solucion}
\end{document}
